\documentclass[11pt,a4paper]{article}

% ---------------------------------------------------------
% PREAMBLE (stable, warning-free, Overleaf-safe)
% ---------------------------------------------------------
\usepackage[utf8]{inputenc}
\usepackage[T1]{fontenc}
\usepackage{lmodern}
\usepackage{amsmath, amssymb, amsfonts}
\usepackage{bm}
\usepackage{geometry}
\usepackage{graphicx}
\usepackage{hyperref}
\usepackage{cite}
\usepackage{authblk}
\usepackage{physics}
\usepackage{mathtools}

\geometry{margin=2.4cm}

\title{\textbf{Companion XXXIX: The Thermal Sunyaev--Zel'dovich Effect in the Relaxation-Driven Cyclic Cosmology}}
\author{Michael Lehmann \\ \texttt{mi.lehmann@gmx.de}}
\date{April 2026}

\begin{document}
\maketitle

\begin{abstract}
The thermal Sunyaev--Zel'dovich (tSZ) effect provides a direct probe of the pressure 
distribution in galaxy clusters and the thermodynamic state of the intracluster 
medium (ICM). In the Relaxation-Driven Cyclic Cosmology (RDCC), the scale-dependent 
evolution of the gravitational potential modifies the formation history of clusters, 
the temperature and density structure of their gas, and the coherence of the 
surrounding large-scale environment. This Companion develops a continuous 
description of the tSZ effect in RDCC, deriving how the relaxation scale influences 
the Compton-$y$ parameter, the pressure profile, and the integrated tSZ signal. The 
resulting framework provides a natural interpretation of observed deviations from 
self-similar scaling relations and offers new predictions for upcoming CMB surveys.
\end{abstract}
% ---------------------------------------------------------
\section{Introduction}

The thermal Sunyaev--Zel'dovich effect arises from the inverse Compton scattering 
of cosmic microwave background photons by hot electrons in galaxy clusters. It is 
a unique probe of the thermal energy content of the intracluster medium and is 
largely insensitive to redshift. In the standard cosmological model, the tSZ signal 
is determined by the depth of the gravitational potential, the virial temperature 
of the gas, and the pressure profile of the cluster.

In RDCC, the gravitational potential evolves differently. The relaxation mechanism 
suppresses the growth of small-scale modes and enhances the coherence of large-scale 
modes. This modifies the collapse history of clusters, the thermodynamic structure 
of their gas, and the pressure distribution that determines the tSZ signal. The 
tSZ effect becomes a direct tracer of the relaxation scale $k_{\mathrm{rel}}$ and 
the temporal structure of the gravitational field.
% ---------------------------------------------------------
\section{The RDCC Potential and Cluster Formation}

The formation of galaxy clusters is governed by the evolution of the gravitational 
potential. In RDCC, this potential takes the form
\begin{equation}
    \Phi_{\mathrm{RDCC}}(k,a)
    = \Phi(k,a)
      \left[1 - \chi_{\Phi}
      \left(\frac{k}{k_{\mathrm{rel}}}\right)^{\alpha}\right].
\end{equation}

Large-scale modes, which dominate the collapse of massive halos, retain a high 
degree of temporal coherence. As a result, clusters form earlier and with deeper 
potential wells than in the standard cosmological model. Small-scale modes, which 
would normally contribute to the steepening of the inner density profile, are 
suppressed, leading to smoother central regions.

The virial temperature of the intracluster medium becomes
\begin{equation}
    T_{\mathrm{vir,RDCC}}
    = T_{\mathrm{vir}}
      \left[1 - \chi_{T}
      \left(\frac{k(M)}{k_{\mathrm{rel}}}\right)^{\alpha}\right],
\end{equation}
reflecting the modified collapse dynamics.
% ---------------------------------------------------------
\section{Pressure Profiles and the Compton-\texorpdfstring{$y$}{y} Parameter}

The tSZ effect is quantified by the Compton-$y$ parameter,
\begin{equation}
    y = \frac{\sigma_{T}}{m_{e} c^{2}}
        \int P_{e}(r)\,dl,
\end{equation}
where $P_{e}(r)$ is the electron pressure.

In RDCC, the pressure profile is modified by the relaxation mechanism. The 
hydrostatic equilibrium condition,
\begin{equation}
    \frac{1}{\rho}\frac{dP}{dr}
    = -\frac{d\Phi_{\mathrm{RDCC}}}{dr},
\end{equation}
yields a pressure distribution that is smoother in the core and more extended in 
the outskirts.

The RDCC pressure profile can be approximated as
\begin{equation}
    P_{\mathrm{RDCC}}(r)
    = P(r)
      \left[
        1 - \chi_{P}
        \left(\frac{1}{k_{\mathrm{rel}} r}\right)^{\alpha}
      \right],
\end{equation}
where $P(r)$ is the standard pressure profile.

This modification enhances the tSZ signal in the outskirts while reducing it in 
the core.
% ---------------------------------------------------------
\section{Integrated tSZ Signal and Scaling Relations}

The integrated tSZ signal is given by
\begin{equation}
    Y = \int y\,d\Omega.
\end{equation}

In RDCC, the integrated signal becomes
\begin{equation}
    Y_{\mathrm{RDCC}}
    = Y
      \left[
        1 - \chi_{Y}
        \left(\frac{k(M)}{k_{\mathrm{rel}}}\right)^{\alpha}
      \right].
\end{equation}

This modification alters the scaling relation between $Y$ and halo mass. Massive 
clusters, whose formation is accelerated by the coherence of large-scale modes, 
exhibit enhanced tSZ signals. Low-mass clusters, whose collapse is delayed by the 
relaxation mechanism, show suppressed tSZ signals.

These deviations from self-similar scaling provide a direct observational probe of 
the relaxation scale.
% ---------------------------------------------------------
\section{Observational Signatures}

The RDCC modifications to the tSZ effect manifest in several measurable ways. 
Clusters exhibit smoother pressure profiles, enhanced tSZ signals in their 
outskirts, and deviations from standard scaling relations. The tSZ power spectrum 
is similarly affected, with a suppression of small-scale power and an enhancement 
of large-scale coherence.

These signatures provide a direct observational window into the relaxation scale 
and the temporal structure of the gravitational field in RDCC.
% ---------------------------------------------------------
\section{Conclusion}

The thermal Sunyaev--Zel'dovich effect in the Relaxation-Driven Cyclic Cosmology 
reflects the scale-dependent evolution of the gravitational potential. The 
relaxation mechanism modifies the formation history of clusters, the thermodynamic 
structure of their gas, and the pressure distribution that determines the tSZ 
signal. These effects produce distinctive signatures in the Compton-$y$ parameter, 
the integrated tSZ signal, and the tSZ power spectrum. The present Companion 
establishes the theoretical foundation for future studies of cluster astrophysics, 
CMB secondary anisotropies, and the interplay between baryons and dark matter in 
RDCC.

% ========================
% References
% ========================

\begin{thebibliography}{10}

\bibitem{Flagship} Lehmann, M. (2026). Relaxation-Driven Cyclic Cosmology (RDCC): A Minimal CPT-Symmetric Framework with a Single Parameter. Flagship Paper. Zenodo: \url{https://zenodo.org/records/19744958} (v43.0 or later).

\bibitem{Zenodo} The complete RDCC companion series (I--LVII) is available at the same Zenodo record above.

% Thematisch relevante Companions
\bibitem{CompXVI} Companion XVI: Boltzmann Code and Modified Hierarchy.
\bibitem{CompXXXI} Companion XXXI: RDCC Halo Model and Non-Linear Structure.
\bibitem{CompXXXII} Companion XXXII--XXXIX: Galaxy--Halo Connection, Cosmic Web, Lensing, RSD, ISW, tSZ, Rotation Curves, CGM.

% Wichtige externe Referenzen
\bibitem{Planck2020} Planck Collaboration (2020), Astron. Astrophys. 641, A6.
\bibitem{DESI2024} DESI Collaboration (2024/2025), arXiv: [aktuelle $f\sigma_8$/BAO-Paper].
\bibitem{JWST2024} Maiolino et al. (2024), Nature (JWST high-$z$ galaxies/quasars).
\bibitem{Kaiser1987} Kaiser, N. (1987), Mon. Not. R. Astron. Soc. 227, 1 (RSD).
\bibitem{Bartelmann2001} Bartelmann, M. \& Schneider, P. (2001), Phys. Rep. 340, 291 (Weak Lensing Review).
\bibitem{HuOkamoto2002} Hu, W. \& Okamoto, T. (2002), Astrophys. J. 574, 566 (CMB Lensing).
\bibitem{LewisChallinor2006} Lewis, A. \& Challinor, A. (2006), Phys. Rep. 429, 1 (CMB Lensing Review).
\bibitem{NFW1997} Navarro, J. F., Frenk, C. S. \& White, S. D. M. (1997), Astrophys. J. 490, 493 (Halo Profiles).

\end{thebibliography}

\end{document}
